\raggedbottom
\chapter{Descripción del problema}
\section{Problema}

La logística de transporte de mercancías es una de las áreas fundamentales para que se pueda desarrollar el comercio tanto a nivel
local como global. El éxito de las operaciones comerciales depende de la eficiencia y seguridad del transporte de mercancías lo que convierte
al transporte de las remesas en un factor importante.\\


\subsection{Problemas generales en las cadenas de suministro}
Según \textcite{litke2019blockchain} en las cadenas de suministro actuales se han identificado los siguientes problemas:
\begin{itemize}
    \item \textbf{Falta de confianza:} Cada parte involucrada en la cadena de suministro necesita confiar en los datos emitidos y compartidos. En las cadenas 
de suministro actuales, este nivel de confianza no existe debido a fallos en el transporte, errores humanos y fraude intencionado.
    \item \textbf{Captura y gestión de datos:} Es esencial que estos sistemas puedan recopilar datos válidos de forma efectiva, registrarlos de forma segura y proporcionar información precisa a otros sistemas, sin embargo estas tareas fallan con frecuencia y povocan ralentizaciones, alteraciones o impedimentos en el flujo de información.
    \item  \textbf{Inconsistencia de la información:} Debido a la falta de estándares altos y de integración extremo a extremo genera inconsistencias en la información entre los distintos actores. Cada parte debe realizar comprobaciones y verificaciones de la validez de los datos recibidos para asegurar la integridad y eficacia, produciendo así ralentizaciones en el sistema.
    \item \textbf{Limitaciones estructurales:} La estructura actual de las cadenas de suministro dificulta la manipulación eficiente del flujo de información. La trazabilidad y el monitoreo de productos son tareas complejas que afectan negativamente al rendimiento y la eficiencia. Por ejemplo, rastrear el sistema completo desde la extracción de materias primas hasta que el producto es entregado al cliente final es un procedimiento muy complejo en las cadenas de suministro modernas.
\end{itemize}

\subsection{El problema de la última milla}
Uno de los grandes desafíos del transporte de remesas es el denominado entrega de última milla.
Según \textcite{problemaUltimaMilla}, este término se refiere al paso final de la cadena de suministro de un producto el cual se caracteriza por trayectos cortos entre un punto de partida central y múltiples puntos finales, o viceversa.
Este término se refiere al paso final de la cadena de suministro de un producto el cual se caracteriza por trayectos cortos entre un punto de partida central y múltiples puntos finales, o viceversa.
Se estima que la parte de la entrega de la última milla como porcentaje del costo total de la entrega es de hasta un 50\% en algunos paises.
Uno de los principales problemas de la última milla a pesar de las innovaciones tecnológicas introducidas hasta la fecha es el "comprobante de entrega". Esta entrega provoca retrasos, errores, ineficiencias y consume mucho tiempo para los implicados.  


\subsection{Experiencia propia}
Por experiencia propia, he notado que cuando compro por internet en plataformas como \textit{Amazon} o \textit{AliExpress} entre otras,
echo en falta una mayor transparencia en el seguimiento de mis pedidos. Un ejemplo claro fue cuando hice un pedido en España de una batería portátil y un adaptador
de enchufe de tipo europeo a británico. Desgraciadamente, cuando llegó el paquete, solo estaba el adaptador; la batería portátil había 
desaparecido en algún punto del proceso. Lo más frustrante fue no tener forma de comprobar en que momento había fallado o desaparecido la batería portátil y encima de todo
la aplicación de amazon me decía que el pedido se había entregado correctamente: ¿se perdió en el proceso de empaquetado, en algún punto intermedio o en el transporte?
Esa falta de información y control hizo que la experiencia fuera realmente frustrante además de ocasionar costes que se pueden evitar a la propia empresa de \textit{Amazon}.

\newpage

\subsection{Opiniones de usuarios en internet}
Si consultamos en Internet, podemos encontrar diversas opiniones negativas sobre el servicio de logística recibido por parte de distintas empresas. Comentarios en redes sociales, foros y plataformas de opinión recogen casos de paquetes que se pierden, entregas fallidas, cambios de estado poco claros o demoras no justificadas.

\begin{enumerate}
    \item En un comentario publicado en el foro de Reddit, disponible en el siguiente enlace: \href{https://www.reddit.com/r/es/comments/1an0c3x/seur_es_una_mierda/?tl=es-es}{Comentario}, un usuario manifiesta su descontento con la empresa de reparto SEUR. El autor del mensaje destaca haber recibido un servicio deficiente en múltiples ocasiones, señalando incidencias como pedidos marcados como entregados que nunca llegaron a su destino, así como entregas realizadas en establecimientos donde no conocía a ninguna persona que pudiera haber recogido el paquete, entre otras situaciones problemáticas.
    \item Al consultar la página \href{https://es.trustpilot.com}{Trustpilot}, una plataforma en la que consumidores y empresas pueden interactuar mediante reseñas, se observa una gran cantidad de valoraciones negativas hacia el servicio de mensajería Correos Express. En concreto, existen más de diez mil opiniones que otorgan una puntuación promedio de tan solo 1,1 sobre 5 estrellas, lo que refleja un elevado nivel de insatisfacción por parte de los usuarios.
\end{enumerate}

Estas experiencias hacen referencia a la etapa final de la cadena de suministro, es decir, la entrega del producto al cliente. Sin embargo, es importante destacar que los retrasos, errores o la falta de transparencia en cualquiera de las etapas previas también repercuten negativamente en la satisfacción del usuario final. La imposibilidad de conocer en qué momento o lugar se ha producido una incidencia genera frustración y desconfianza tanto en los consumidores como en las propias empresas.

En conjunto, estos testimonios refuerzan la percepción de que la cadena logística actual presenta carencias significativas en términos de trazabilidad, transparencia y responsabilidad.


% \section{Solución propuesta}
% \subsection{Definición de Blockchain}
% 







    
    
    